% A.76: density-one rigidity of completed multiplicative classes.
\begin{statement}{A.76}{Proved}{Density-one detection of completed Kummer classes}
Let $\mathcal P$ be a set of rational primes of Dirichlet
density one.  The continuous localization homomorphism
\[
       \widehat{\Q^\times}\longrightarrow
                      \prod_{p\in\mathcal P}\widehat{\Q_p^\times}
       \tag{A.76.1}
\]
is injective.  More precisely, every nontrivial completed
class is nontrivial locally at every prime in some nonempty
Chebotarev conjugacy class, apart from finitely many primes.

Let $S\subset\Q$ be nonempty and finite, let $R\in\Q\setminus S$,
and suppose that ordinary local coordinates
$x_p\in\Q_p\setminus S$ and global classes
$U_c\in\widehat{\Q^\times}$ satisfy
\[
             (U_c)_p=\kappa_p(x_p-c)
                          \quad(c\in S,\ p\text{ finite}).
       \tag{A.76.2}
\]
If $x_p=R$ on a density-one set of finite primes, then
\[
          U_c=R-c\text{ in }\widehat{\Q^\times}\quad(c\in S),
          \qquad x_p=R\quad\text{at every finite prime}.
       \tag{A.76.3}
\]
If an ordinary real coordinate also satisfies the completed
comparisons, it lies in the same component of $\R\setminus S$
as $R$.  Equality of the real numbers is not required by these
finite-puncture comparisons.
\end{statement}
\begin{proof}
We first prove the finite-level input with its precise loss
of exponent.  Let $m\ge1$ and $a\in\Q^\times$, and suppose
$a$ is a $2m$th power in $\Q_p$ for every $p$ in a set of
Dirichlet density one.  Put
\[
      F=\Q(\mu_{2m}),\qquad
      L=F\bigl(a^{1/(2m)}\bigr),\qquad
      N=\operatorname{Gal}(L/F).
       \tag{A.76.4}
\]
The field $L$ is the splitting field of $Z^{2m}-a$ over
$\Q$.  At an unramified prime away from $2m$ and the numerator
and denominator of $a$, local solubility of this polynomial
is equivalent to its Frobenius permutation having a fixed
root.  Indeed, the roots are integral units, reduction is
separable, and a Frobenius-fixed residue root lifts uniquely
by Hensel's lemma; conversely a local root is fixed by the
local Frobenius action in the unramified splitting field.

If $1\ne\sigma\in N$, it multiplies a chosen root by a
nontrivial $2m$th root of unity and fixes every root of unity.
It therefore fixes none of the $2m$ roots.  Neither does any
conjugate of $\sigma$.  Chebotarev gives a positive-density
set of primes in this conjugacy class
\cite[Theorem 8.31]{AgentMilneANT}.  Removing a density-zero
set and finitely many primes cannot remove that set, a
contradiction.  Hence $N=1$ and every root lies in $F$.

Also $a>0$.  If $a<0$, no root of $Z^{2m}-a$ is real,
so complex conjugation in $\operatorname{Gal}(L/\Q)$ fixes
no root.  Its conjugacy class gives the same Chebotarev
contradiction.  Let $\alpha=a^{1/(2m)}>0$ be the positive
real root.  It lies in $F$, because all roots lie there.
The field $\Q(\alpha)$ is a subfield of the abelian Galois
extension $F/\Q$, so is itself Galois over $\Q$.
It is contained in $\R$, and all its conjugates of $\alpha$
must be real roots of the same polynomial, hence must belong
to $\{\alpha,-\alpha\}$.  It follows that $\alpha^2\in\Q$:
this is immediate in degree one, and in degree two the
nontrivial conjugate is $-\alpha$.
Consequently
\[
                        a=(\alpha^2)^m\in(\Q^\times)^m.
       \tag{A.76.5}
\]
This proves the finite-level assertion with exponent $2m$
in the hypothesis and exponent $m$ in the conclusion.

Now let $U\in\widehat{\Q^\times}$ be in the kernel of
(A.76.1).  For each $m$ choose a rational representative
$a_{2m}$ of its coordinate in
$\Q^\times/(\Q^\times)^{2m}$.  Localization and reduction
modulo $2m$ commute, so $a_{2m}$ is a $2m$th power at
every prime of $\mathcal P$.  By (A.76.5) its image modulo
$m$ is trivial.  Compatibility of the inverse-limit
coordinates says that this image is exactly $U\bmod m$.
It is trivial for every $m$, hence $U=1$.

For the more precise detection statement, choose $m$ with
$U\bmod m\ne1$ and form (A.76.4) from a representative
$a_{2m}$.  If $N\ne1$, the conjugacy class of a nonidentity
element of $N$ has no fixed root.  If $N=1$ and $a_{2m}<0$,
use complex conjugation.  The only remaining case,
$N=1$ and $a_{2m}>0$, would give (A.76.5), contradicting
$U\bmod m\ne1$.  Thus one of the first two explicit
conjugacy classes always exists.  At its good unramified
primes $U\bmod2m$ is nontrivial locally, so $U$ itself is
nontrivial locally.  Chebotarev gives the claimed positive
density.

For (A.76.2), the class $U_c/(R-c)$ is trivial locally at
every prime where $x_p=R$.  Injectivity of (A.76.1) gives
$U_c=R-c$.  At an arbitrary finite prime, the equality
\[
                     \kappa_p(x_p-c)=\kappa_p(R-c)
       \tag{A.76.6}
\]
then implies $x_p=R$, since
$\Q_p^\times\to\widehat{\Q_p^\times}$ is injective by A.2.
One fixed $c\in S$ already suffices for this step.
At the real place, A.2 instead identifies the completed
group with the sign group.  Its comparisons are exactly
\[
            \operatorname{sgn}(x_\infty-c)
                     =\operatorname{sgn}(R-c)\quad(c\in S),
       \tag{A.76.7}
\]
which characterize the component of $\R\setminus S$
containing $R$.
\end{proof}

\paragraph{Why the doubled exponent is necessary.}
The corresponding assertion at one fixed exponent is false.
The number $16$ is an eighth power in $\Q_p$ for every odd
prime $p$, and in $\R$, but is not an eighth power in
$\Q$ or $\Q_2$.  To prove the local assertion, at least
one of $2,-2,-1$ is a square in $\Q_p$: their quadratic
characters have product one, so cannot all be negative.
The respective eighth roots of $16$ are
\[
                       \sqrt2,\qquad\sqrt{-2},\qquad1+i.
       \tag{A.76.8}
\]
The last identity uses $(1+i)^8=16$; the first two are
immediate.  The unit-square criterion and simple-root
Hensel lifting justify the square roots for every odd $p$.
At two, the valuation $v_2(16)=4$ is not divisible by eight.
This is the precise finite-level phenomenon avoided by
passing from level $2m$ to level $m$ in (A.76.5).
No claim of finite-level injectivity is used.

\paragraph{Application to the fixed tuples of A.69 and A.73.}
The tuple in A.69 equals its rational base $R=1/17$ at every
finite prime except two.  The tuple in A.73 equals its base
$R=0$ at every finite prime except two.  Both sets have
density one.  If their old finite-level classes could be
replaced by global completed old classes with these exact
local coordinates, (A.76.3) would force their two-adic
coordinate to be their base as well.  This contradicts the
displayed two-adic values $16/17$ and $1$, respectively.
Thus these particular tuples already fail to have completed
old coordinates.  Their new-coordinate obstructions remain
valid finite-level counterexamples, with the exact missing
passage now proved: a completed example differing from a
rational point must differ at a set of finite primes
containing a positive-density Chebotarev subset.

The certificate \path{code/density_one_kummer_certificate.py}
checks the finite Galois root action and factorization for
(A.76.8), together with the odd-prime square-cover tests.
The density-one theorem itself is proved by the argument above.
